\documentclass[11pt,a4paper]{article}
\usepackage[utf8]{inputenc}
\usepackage[margin=2.4cm]{geometry}
\usepackage{booktabs,enumitem,longtable}
\title{SafeSlice, Step 2: Terminology\\ \large All terms of the domain, around 105 in total}
\author{SafeSlice group}
\date{v1.0, freeze on Friday}
\begin{document}
\maketitle

\noindent Naming convention for everything: CamelCase, singular, plain ASCII, and no parent name inside a class name (the hierarchy already carries that). Named pizzas keep the word Pizza only to avoid clashes with toppings (PepperoniPizza versus the topping Pepperoni).

\section{Classes by branch (72)}

\subsection*{Pizza (11)}
MargheritaPizza, PepperoniPizza, QuattroFormaggiPizza, HawaiianPizza, VegetarianaPizza, PestoSpecialePizza, TonnoPizza, VeganGardenPizza, DiavolaPizza, FunghiPizza

\subsection*{Ingredient (37)}
\begin{itemize}[leftmargin=*, itemsep=1pt]
\item \textbf{Vegetable}: Tomato, Onion, Olive, Rucola, Garlic; \textbf{Herb}: Basil, Oregano; \textbf{PepperVegetable}: BellPepper, Jalapeno
\item \textbf{Cheese}: Mozzarella, Parmesan, Gorgonzola, Provola, VeganCheese
\item \textbf{Meat}: Ham, Salami, Pepperoni, Chicken, PorkGelatin
\item \textbf{Fish}: Tuna, Anchovy
\item \textbf{Fruit}: Pineapple
\item \textbf{Fungus}: Mushroom
\item \textbf{Nut}: PineNut
\item \textbf{Sauce}: TomatoSauce, Pesto, GarlicOil
\item \textbf{Spice}: BlackPepper, ChiliFlakes
\item \textbf{BaseIngredient}: WheatFlour, GlutenFreeFlour, Water, Yeast, Salt, OliveOil, Egg
\item \textbf{Component}: Gluten, Milk, Pork
\end{itemize}

\subsection*{DoughBase (4)}
NormalDough, GlutenFreeDough, VeganDough

\subsection*{DietaryRestriction (11)}
\textbf{Diet}: VeganDiet, VegetarianDiet, PescatarianDiet, HalalDiet, KosherDiet.\\
\textbf{Allergy}: NutAllergy, GlutenIntolerance, LactoseIntolerance.\\
Note: an allergy is not a diet, so they are two separate children of DietaryRestriction.

\subsection*{Customer (1), Certification (4), SpicinessLevel (4)}
Customer. HalalCertified, KosherCertified, VeganCertified. NotSpicy, MediumSpicy, HotSpicy (a value partition: disjoint and covering).

\section{Properties (11)}
\begin{tabular}{lll}
\toprule
Property & Signature & Typed or derived\\
\midrule
hasTopping & Pizza to Ingredient & typed\\
hasBase & Pizza to DoughBase & typed\\
madeFrom & Ingredient to Ingredient & typed\\
contains & Pizza to Ingredient & derived\\
hasSpiciness & Ingredient to SpicinessLevel & typed\\
hasRestriction & Customer to DietaryRestriction & typed\\
forbids & DietaryRestriction to category & typed\\
hasCertification & Ingredient to Certification & typed\\
unsafeFor & Pizza to Customer & derived\\
safeFor & Pizza to Customer & derived\\
\bottomrule
\end{tabular}

\section{Defined classes (10)}
VegetarianPizza, VeganPizza, PescatarianPizza, HalalPizza, NutFreePizza, GlutenFreePizza, DairyFreePizza, SpicyPizza, MeatLoversPizza, and VegetarianSalamiPizza (deliberately impossible, our contradiction demo).

\section{Individuals (12)}
Customers: anna (VeganDiet), omar (HalalDiet and NutAllergy), mia (LactoseIntolerance), lisa (VegetarianDiet and GlutenIntolerance), tom (no restrictions). Plus example pizza individuals if the OWL team wants an ABox demo.

\section{Count}
72 classes + 11 properties + 10 defined classes + 12 individuals, roughly 105 terms, inside the 50 to 100 target with margin.
\end{document}
